\documentclass{article}

\usepackage{changepage}
\usepackage{xcolor}
\usepackage{amsmath}
\usepackage{amssymb}
\usepackage{tikz}
\usepackage{tcolorbox}

\definecolor{blue}{RGB}{0, 166, 223}
\definecolor{lightblue}{RGB}{220, 240, 247}

\newcommand{\theorem}[2]{
	\vspace{0.25cm}
	\begin{tcolorbox}[
		colback=lightblue,
		colframe=blue,
		boxrule=0.5pt,
		toprule=2pt,
		arc=0pt,
		left=0.4cm,
		right=0.4cm,
		top=0.25cm,
		bottom=0.35cm
	]
		\textcolor{blue}{\textbf{Theorem #1}}

		\vspace{0.15cm}

		#2
	\end{tcolorbox}
	\vspace{0.25cm}
}

\newcommand{\header}[1]{
	\vspace{0.25cm}
		\begin{adjustwidth}{0cm}{0cm}
			\textcolor{blue}{\textbf{#1}}
		\end{adjustwidth}
	\vspace{0.25cm}
}

\newcommand{\question}[2]{
	\vspace{0.125cm}
		\begin{adjustwidth}{0.25cm}{0cm}
			\textbf{#1} #2
		\end{adjustwidth}
	\vspace{0.125cm}
}

\newcommand{\container}[1]{
	\vspace{0.25cm}
		\begin{adjustwidth}{0cm}{0cm}
			#1
		\end{adjustwidth}
	\vspace{0.25cm}
}

\newcommand{\solution}[1]{
	\vspace{0.25cm}
		\begin{adjustwidth}{0cm}{0cm}
			\textbf{Solution:} #1
		\end{adjustwidth}
	\vspace{0.25cm}
}

\newcommand{\wrapper}[1]{
	\vspace{0.25cm}
		\begin{adjustwidth}{0cm}{0cm}
			#1
		\end{adjustwidth}
	\vspace{0cm}
}

\newcommand{\calc}[1]{
	\[
		#1
	\]
}

\begin{document}

	\theorem{5.2.2 Sum of a Geometric Sequence}{
		For any real number r except 1, and any integer
		\(n \geq 0\),

		\calc{\sum_{i=0}^{n} r^{i} = \frac{r^{n+1} - 1}{r - 1}}
	}

	\container {
		\textcolor{blue}{\textbf{Proof (by mathematical induction):}}
		Suppose r is a particular but arbitrarily chosen real
		number that is not equal to 1, and let the property P(n) be
		the equation

		\calc{\sum_{i=0}^{n} r^{i} = \frac{r^{n+1} - 1}{r - 1}}

		\wrapper{
			We must show that P(n) is true for every integer
			\(n \geq 0\). We do this by mathematical
			induction on \emph{n}.
		}
	}

	\container{
		\textbf{Show that P(0) is true:}

		\wrapper{
			To etablish \emph{P(0)}, we must show that
		}

		\calc{\sum_{i=0}^{0} r^{i} = \frac{r^{0+1} - 1}{r - 1}}

		\wrapper{
			The left hand side of this equation is \(r^{0}\)
			and the right-hand side is
		}

		\calc{\frac{r^{0+1} - 1}{r - 1} = \frac{r - 1}{r - 1} = 1}

		\wrapper{
			also because \(r^{1} = r\) and, since \(r \neq 1, r - 1 \neq 0\).
			Hence P(0) is true.
		}
	}

	\container{
		\textbf{Show that for every integer \(k \geq 1\), if
		\(P(k)\) is true then \(P(k + 1)\) is also true:}

		\wrapper{
			\emph{[Suppose that \emph{P(k)} is true for a particular
			but arbitrarily chosen integer \(k \geq 1\). That is:]}
			Let \emph{k} be any integer with \(k \geq 0\), and suppose that
		}

		\calc{\sum_{i=0}^{k} r^{i} = \frac{r^{k+1} - 1}{r - 1}}

		\wrapper{
			\emph{[We must show that \emph{P(k + 1)} is true. That is:]}
			We must show that
		}

		\calc{\sum_{i=0}^{k+1} r^{i} = \frac{r^{(k+1)+1} - 1}{r - 1}}

		\wrapper{or, equivalently, that}

		\calc{\sum_{i=0}^{k+1} r^{i} = \frac{r^{k+2} - 1}{r - 1}}

		\wrapper{
			\emph{[We will show that the left-hand side of P(k + 1)
			equals the right-hand side.]}
		}

		\wrapper{the left-hand side of \(P(k + 1)\) is}

		\calc{\sum_{i=0}^{k+1} r^{i} = \sum_{i=0}^{k} r^{i} + r^{k+1}}

		\calc{
			\begin{aligned}
				\sum_{i=0}^{k+1} r^{i} &= \sum_{i=0}^{k} r^{i} + r^{k+1} \\
				&= \frac{r^{k+1} - 1}{r - 1} + r^{k+1}\\
				&= \frac{r^{k+1} - 1}{r - 1} + \frac{r^{k+1}(r - 1)}{r - 1}\\
				&= \frac{(r^{k+1} - 1) + r^{k+1}(r - 1)}{r - 1}\\
				&= \frac{r^{k+1} - 1 + r^{k+2} - r^{k+1}}{r - 1}\\
				&= \frac{r^{k+2} - 1}{r - 1}
			\end{aligned}
		}

		\wrapper{Which is the right-hand side of P(k + 1) \emph{[as was to be shown.]}}
	}

	\container{
		\emph{[Since we have proved the basis step and the inductive step,
		we conclude that the theorem is true.]}
	}


\end{document}
